%-----------------------------------------------------------------------------------------------------------------------------------------------%
%	The MIT License (MIT)
%
%	Copyright (c) 2021 Jitin Nair
%
%	Permission is hereby granted, free of charge, to any person obtaining a copy
%	of this software and associated documentation files (the "Software"), to deal
%	in the Software without restriction, including without limitation the rights
%	to use, copy, modify, merge, publish, distribute, sublicense, and/or sell
%	copies of the Software, and to permit persons to whom the Software is
%	furnished to do so, subject to the following conditions:
%	
%	THE SOFTWARE IS PROVIDED "AS IS", WITHOUT WARRANTY OF ANY KIND, EXPRESS OR
%	IMPLIED, INCLUDING BUT NOT LIMITED TO THE WARRANTIES OF MERCHANTABILITY,
%	FITNESS FOR A PARTICULAR PURPOSE AND NONINFRINGEMENT. IN NO EVENT SHALL THE
%	AUTHORS OR COPYRIGHT HOLDERS BE LIABLE FOR ANY CLAIM, DAMAGES OR OTHER
%	LIABILITY, WHETHER IN AN ACTION OF CONTRACT, TORT OR OTHERWISE, ARISING FROM,
%	OUT OF OR IN CONNECTION WITH THE SOFTWARE OR THE USE OR OTHER DEALINGS IN
%	THE SOFTWARE.
%	
%
%-----------------------------------------------------------------------------------------------------------------------------------------------%

%----------------------------------------------------------------------------------------
%	DOCUMENT DEFINITION
%----------------------------------------------------------------------------------------

\documentclass[a4paper,12pt]{article}

\usepackage{ctex} 

%----------------------------------------------------------------------------------------
%	PACKAGES
%----------------------------------------------------------------------------------------
\usepackage{url}
\usepackage{parskip} 	

%other packages for formatting
\RequirePackage{color}
\RequirePackage{graphicx}
\usepackage[usenames,dvipsnames]{xcolor}
\usepackage[scale=0.9]{geometry}

%tabularx environment
\usepackage{tabularx}

%for lists within experience section
\usepackage{enumitem}

% centered version of 'X' col. type
\newcolumntype{C}{>{\centering\arraybackslash}X} 

%to prevent spillover of tabular into next pages
\usepackage{supertabular}
\usepackage{tabularx}
\newlength{\fullcollw}
\setlength{\fullcollw}{0.47\textwidth}

%custom \section
\usepackage{titlesec}				
\usepackage{multicol}
\usepackage{multirow}

%CV Sections inspired by: 
%http://stefano.italians.nl/archives/26
\titleformat{\section}{\Large\scshape\raggedright}{}{0em}{}[\titlerule]
\titlespacing{\section}{0pt}{10pt}{10pt}

%for publications
\usepackage[style=authoryear,sorting=ynt, maxbibnames=2]{biblatex}

%Setup hyperref package, and colours for links
\usepackage[unicode, draft=false]{hyperref}
\definecolor{linkcolour}{rgb}{0,0.2,0.6}
\hypersetup{colorlinks,breaklinks,urlcolor=linkcolour,linkcolor=linkcolour}
\addbibresource{citations.bib}
\setlength\bibitemsep{1em}

%for social icons
\usepackage{fontawesome5}

%debug page outer frames
%\usepackage{showframe}


% job listing environments
\newenvironment{jobshort}[2]
    {
    \begin{tabularx}{\linewidth}{@{}l X r@{}}
    \textbf{#1} & \hfill &  #2 \\[3.75pt]
    \end{tabularx}
    }
    {
    }

\newenvironment{joblong}[2]
    {
    \begin{tabularx}{\linewidth}{@{}l X r@{}}
    \textbf{#1} & \hfill &  #2 \\[3.75pt]
    \end{tabularx}
    \begin{minipage}[t]{\linewidth}
    \begin{itemize}[nosep,after=\strut, leftmargin=1em, itemsep=3pt,label=--]
    }
    {
    \end{itemize}
    \end{minipage}    
    }

%----------------------------------------------------------------------------------------
%	BEGIN DOCUMENT
%----------------------------------------------------------------------------------------

\begin{document}

\pagestyle{empty} 

%----------------------------------------------------------------------------------------
%	TITLE
%----------------------------------------------------------------------------------------

\begin{tabularx}{\linewidth}{@{} C @{}}
\Huge{左天佑} \\[7.5pt]
\href{mailto:fztym@sjtu.edu.cn}{\raisebox{-0.05\height}\faEnvelope \ fztym@sjtu.edu.cn} \ $|$ \ 
\href{tel:+8613162909941}{\raisebox{-0.05\height}\faMobile \ +86.131.6290.9941} \\
\end{tabularx}

%----------------------------------------------------------------------------------------
%	EDUCATION
%----------------------------------------------------------------------------------------

\section{教育背景}
\begin{tabularx}{\linewidth}{@{}l X@{}}	
\textbf{上海交通大学} & 中国上海 \\
\textit{电子计算机工程，计算机科学辅修} & 2023 年 9 月 - 至今 \\  
\normalsize (GPA: 3.62/4.0，排名 51/215) & 核心专业课均获 4.0/4.0\\
\end{tabularx}

%----------------------------------------------------------------------------------------
%	RESEARCH INTERESTS
%----------------------------------------------------------------------------------------

\section{研究兴趣}
高性能计算与 GPU 体系结构、算法设计与分析。

%----------------------------------------------------------------------------------------
%	PUBLICATIONS
%----------------------------------------------------------------------------------------

\section{学术论文}
\begin{tabularx}{\linewidth}{@{}X@{}}
\textbf{Disproofs of the Four-Coin Hereditary Subcurrency Conjecture and the Prefix Pattern Representative Conjecture.} \normalsize（在投 SODA 2027）研究换零钱问题中贪心算法保持最优的 orderly 货币系统，构造反例推翻了关于遗传性子货币模式的两项已有猜想。 \\
\end{tabularx}

%----------------------------------------------------------------------------------------
%	RESEARCH & ENGINEERING EXPERIENCE
%----------------------------------------------------------------------------------------

\section{科研与工程经历}

\begin{joblong}{阿里巴巴 | 推理引擎 · 算子优化实习生}{2026 年 7 月 - 至今}
\item 参与大模型推理引擎的算子优化，围绕 \textbf{Qwen3-Next} 中的线性注意力模块开展工作，涉及 \textbf{Gated DeltaNet (GDN)} 与 \textbf{fused-recurrent} 算子。
\end{joblong}

\begin{joblong}{摩尔线程 | GPU 计算与 AI 软件实习生}{2026 年 1 月 - 2026 年 6 月}
\item 参与自研高性能 GPU 数学库 muBLAS（对标 NVIDIA cuBLAS）的核心算子研发，独立负责批处理稠密线性代数例程：LU 分解 \texttt{getrfBatched}、线性方程组求解 \texttt{getrsBatched}、矩阵求逆 \texttt{getriBatched} / \texttt{matinvBatched} 的设计与实现。
\item 系统分析共享内存布局、访存合并与归约策略对性能与数值稳定性的影响，通过 profiling 定位瓶颈并迭代优化内核，形成“建模—实验—验证”的完整性能调优方法论。
\item 调研 cuBLAS / MAGMA 等业界实现与相关文献，对比不同选主元（pivoting）策略在批处理小矩阵场景下的精度—性能权衡。
\end{joblong}

\begin{joblong}{LemonDB：现代 C++ 多线程内存数据库}{\href{https://github.com/zty2004/ECE4821Jp2}{项目链接}}
\item 基于 \textbf{C++20} 实现高并发关系型数据库，设计拉取式线程池（批量拉取任务将全局锁竞争降低 \textbf{16 倍}）与依赖感知调度系统，自动解析查询间资源依赖并进行拓扑调度。
\end{joblong}

\begin{joblong}{MSD 分布式音乐推荐系统}{\href{https://github.com/zty2004/ECE4721J}{项目链接}}
\item 基于百万歌曲数据集构建端到端大数据管道（288GB 原始数据转为 884MB Avro，PySpark 处理耗时从 5.5 小时缩短至 \textbf{1.3 小时}），并用 FAISS 实现亚秒级近似最近邻推荐。
\end{joblong}

%----------------------------------------------------------------------------------------
%	TEACHING
%----------------------------------------------------------------------------------------

\section{教学经历}

\begin{jobshort}{ECE4770J 算法导论 | 助教}{2025 年 9 月 - 2025 年 12 月}
开展 OCaml 算法设计与实现实验课，重点讲解尾递归优化与延续传递风格 (CPS)；维护并扩展在线判题系统 (OJ)，强化测试用例并调试评分脚本。
\end{jobshort}

\begin{jobshort}{ECE2810J 数据结构与算法 | 助教}{2025 年 9 月 - 2025 年 12 月}
主持习题课，讲授高级数据结构与复杂度分析；设计受信息学竞赛启发的补充习题集，并维护课程 OJ 系统。
\end{jobshort}

%----------------------------------------------------------------------------------------
%	Courses
%----------------------------------------------------------------------------------------

\section{主修课程}
\begin{tabularx}{\linewidth}{@{}l X@{}}	
ECE4770J & 算法导论 \hfill \normalsize (4.0/4.0)\\
ECE4821J & 操作系统导论 \hfill \normalsize (4.0/4.0)\\
ECE4720J & 大数据方法与工具 \hfill \normalsize (4.0/4.0)\\
ECE2810J & 数据结构与算法 \hfill \normalsize (4.0/4.0)\\
MATH1860J & 荣誉数学 II \hfill \normalsize (4.0/4.0)\\
ENGR1510J & 计算机导论（进阶） \hfill \normalsize (4.0/4.0)
\end{tabularx}

%----------------------------------------------------------------------------------------
%	Prizes & Honors
%----------------------------------------------------------------------------------------

\section{奖项和荣誉}
\begin{tabularx}{\linewidth}{@{}l X@{}}
\textbf{一等奖}, & NOIP 2020（全国青少年信息学奥林匹克联赛） \\
\end{tabularx}

%----------------------------------------------------------------------------------------
%	SKILLS
%----------------------------------------------------------------------------------------

\section{技能专长}
\begin{tabularx}{\linewidth}{@{}l X@{}}
编程语言 &  \normalsize{C++, Python, OCaml, Shell} \\
系统与工具 & \normalsize{CUDA, Linux, Git, Docker, vLLM, Spark} \\
语言能力 &  \normalsize{普通话（母语），英语（全英文工科培养体系）} \\
\end{tabularx}

\vfill
\center{\footnotesize 最后更新：\today}

\end{document}
